% Ad Familiares 9.12 --- headnote
% ad-familiares-09-12, late 45 BC, Cumanum

\textit{Cicero to P. Cornelius Dolabella, written at
Cicero's Cumanum villa in (probably) late 45 BC. The
Perseus dateline is partially abbreviated --- \textit{in
Cumano mcd. m. Dcc. a.~709 (45)} --- and the meta entry
carries the year-precision placeholder \textit{-0045-01-17};
the strongest internal indicator is the mention of the
\textit{oratiuncula pro Deiotaro}, the little speech for
King Deiotarus, which Cicero delivered before Caesar in
Rome in November 45 and which here he sends to Dolabella
in writing. The letter is therefore late autumn or winter
of 45.}

\textit{Two short paragraphs, each in a different
register. The first is a playful tease: Cicero
congratulates the people of Baiae on the resort's having
suddenly become healthy, then suggests, mock-suspicious,
that perhaps Baiae has only put on this charm to please
Dolabella, and that sky and earth themselves will set
aside their nature if it suits him. The second turns to
the speech and to a personal injury Dolabella has
suffered (unspecified, but possibly the matrimonial
strain after his divorce from Tullia, or some political
slight): Cicero downplays the speech with the lovely
self-deprecating phrase \textit{levidense crasso filo} ---
``a little gift of slight stuff and coarse thread''
--- and closes with the recommendation that Dolabella
bear his troubles wisely, so that his own moderation may
shame the injustice of others. \textbf{Metadata note: the
\texttt{meta/works.yaml} entry carries the placeholder
day \texttt{-0045-01-17} at year precision, which
disagrees with the Perseus dateline's location of the
letter at \textit{Cumano} in 45; the internal reference to
the Pro Deiotaro points to late 45, after the speech's
November delivery. The entry's date should be revised
when the metadata is consolidated.}}
